\documentclass{report}
\usepackage[utf8]{inputenc}
\usepackage[spanish, es-noshorthands]{babel}
\usepackage{amsmath}
\usepackage{amssymb}

\begin{document}

\chapter*{Resumen}

El problema de la satisfacibilidad booleana (SAT) se estudia en este
trabajo desde la teoría de lenguajes formales. En la línea de
investigación en la que se inscribe, la satisfacibilidad de una fórmula
se reduce al vacío de la intersección de un lenguaje regular de
interpretaciones con una gramática que impone la consistencia entre las
ocurrencias de cada variable booleana.
Con gramáticas menos expresivas esa reducción se limita a fórmulas con
cierto orden de las ocurrencias: una gramática libre del contexto no
admite cruces, y en una de concatenación de rango simple (sRCG) la aridad
crece con la fórmula. En un intento de resolver más fórmulas en tiempo
polinomial se ataca el problema con gramáticas de concatenación de rango
(RCG): con su no linealidad se modela la consistencia de más fórmulas y
se mantiene mínima la aridad, como se muestra con una RCG que reconoce la
consistencia de cualquier fórmula de SAT con aridad uno; pero por esa
misma no linealidad se pierde la garantía directa de que el vacío sea
polinomial.
Se da un algoritmo de fuerza bruta que decide ese vacío: se enumeran las
cadenas del autómata y se reconocen con la RCG; su costo es exponencial,
así que
no es una vía factible, pero muestra por qué conviene estudiar el
problema sobre un único objeto. Se construye ese objeto, la gramática
producto, y se caracterizan su tamaño y el costo de construirlo; sobre él
se puede abordar el vacío: hallar un método polinomial para todas las
fórmulas, lo que implicaría $P = NP$; probar que no existe; o hallarlo al
restringir a una subclase. Aunque su
tamaño no es polinomial, queda como objeto de estudio: eliminar sus
cláusulas no productivas y analizar el costo de decidir el vacío se
proponen como trabajo futuro.

\vspace{1em}
\noindent\textbf{Palabras clave:} satisfacibilidad booleana, gramáticas
de concatenación de rango, vacío de la intersección, lenguaje regular
finito, no linealidad.

\end{document}
